% SEGMENT OLP-0654-S01
% Part: normal-modal-logic
% Chapter: syntax-and-semantics
% Section: normal-modal-logics

\documentclass[../../../include/open-logic-section]{subfiles}

\begin{document}

\olfileid{nml}{syn}{nor}

\olsection{இயல்பான வாய்ப்புநிலைத் தருக்கங்கள்}

தகுந்த ஒழுங்குப் பண்புகள் உள்ள !!{formula}s இன் எந்தக் கணத்தையும்
வாய்ப்புநிலைத் தருக்கமாகக் கருதலாம். சில வாய்ப்புநிலைத் தருக்கங்கள்
அவற்றின் சுவாரஸ்யமான பொருள்கோள்களாலும் பயன்பாடுகளாலும்,
முக்கியமாக அவற்றைப் பற்றிய கோட்பாட்டு அறிவு நன்கு வளர்ந்திருப்பதாலும்,
சிறப்பாகக் கவனிக்கப்படுகின்றன. இவையே இயல்பான வாய்ப்புநிலைத்
தருக்கங்கள் எனப்படுகின்றன.

% SEGMENT OLP-0654-S02
\begin{defn}\ollabel{defn:modal-logic}
\emph{வாய்ப்புநிலைத் தருக்கம்}~$\Log L$ என்பது அடிப்படை
வாய்ப்புநிலை மொழியின் !!{formula}s இன் கணம்; அது
\begin{enumerate}
\item கூற்றுத் தருக்க மெய்மங்கள் அனைத்தையும் கொண்டிருக்கும்;
\item சீரான பதிலீட்டின் கீழ் மூடப்பட்டிருக்கும்; மேலும்
\item மோடஸ் போனென்ஸின் கீழ் மூடப்பட்டிருக்கும்; அதாவது
  $!A$ மற்றும் $(!A \lif
  !B) \in \Log L$ எனில், $!B \in \Log L$ உம் பொருந்தும்.
\end{enumerate}
\end{defn}

% SEGMENT OLP-0654-S03
\begin{ex}
இவ்வரையறையை நிறைவுசெய்வதால் வாய்ப்புநிலைத் தருக்கங்களாகக்
கணிக்கப்படும், பெரிதாக ஆர்வமூட்டாத !!{formula}s இன்
கணங்களுக்குச் சில எடுத்துக்காட்டுகள்:
\begin{enumerate}
\item மெய்மங்கள் அனைத்தின் கணம்;
\item எல்லா !!{formula}s இன் கணம் (முரணுடைய வாய்ப்புநிலைத் தருக்கம்).
\end{enumerate}
\end{ex}

% SEGMENT OLP-0654-S04
\begin{defn}\label{defn:normal-modal-logic}
வாய்ப்புநிலைத் தருக்கம்~$\Log L$ \emph{இயல்பானது}
எனில், எனில் மட்டுமே,
\begin{enumerate}
\item அது பின்வரும் இரண்டையும் கொண்டிருக்கும்:
\begin{align}
\tag{K} \Box(p \lif q) & \lif (\Box p \lif \Box q) \\
\tag{Dual} \Diamond p & \liff \lnot \Box \lnot p
\end{align}
\item மேலும் அது இன்றியமையாக்கலின் கீழ் மூடப்பட்டிருக்கும்;
  அதாவது $!A \in
  \Log L$ எனும்போதெல்லாம் $ \Box !A \in \Log L$ உம் பொருந்தும்.
\end{enumerate}
\end{defn}

% SEGMENT OLP-0654-S05
வாய்ப்புநிலைத் தருக்கங்களை வரையறுக்கும் இரண்டு முதன்மை வழிகளை
நாம் காண்போம். ஒன்று நிறுவல்-கோட்பாட்டு வழி: குறிப்பிட்ட
!!{derivation} முறைமைகளில் வருவிக்கத்தக்க !!{formula}s இன்
கணங்களாக அவற்றைக் கொள்வது. மற்றொன்று மாதிரிக் கோட்பாட்டு வழி:
குறிப்பிட்ட மாதிரி வகுப்புகளில் செல்லுபடியாகும் !!{formula}s
இன் கணங்களாக அவற்றைக் கொள்வது. இது சாதாரண கூற்றுத்
தருக்கத்திலும் (அதேபோல் முதல்தரத் தருக்கத்திலும்) நிலவும்
நிலையைப் பிரதிபலிக்கிறது. அங்கும் ஓர் !!{formula}s இன்
கணத்தை இரு வேறு வழிகளில் குறிப்பிடலாம்: எடுத்துக்காட்டாக,
எல்லா மெய்மதிப்பு ஒதுக்கீடுகளிலும் மெய்யான மெய்மங்களின்
கணமாகவும், ஏதோவொரு !!{derivation} முறைமையில்
!!{derivable} ஆகும் எல்லா !!{formula}s இன்
கணமாகவும் குறிப்பிடலாம். இயல்பான வாய்ப்புநிலைத் தருக்கங்களுக்கு,
தொடர்புசார் மாதிரிகளையும் அடிகோள்சார் (ஹில்பர்ட் வகை)
நிறுவல் முறைமைகளையும் பயன்படுத்தி இந்த இருவழி
குறிப்பீட்டை அமைக்கலாம்.
\end{document}
